% TOP_PROBLEM: {"id": 3153, "title": "Complete bipartite subgraphs of perfect graphs", "category_id": 3, "category": {"id": 3, "name": "graph_theory", "display_name": "Graph Theory", "description": "Problems involving graphs, networks, and their properties.", "slug": "graph-theory", "order_index": 3, "created_at": "2026-07-31T15:26:25.671Z"}, "status": "open"}
% FIRST_POSTED: null
% ENTRY_KIND: statement_only
% Imported from: batch02.tex; UnsolvedMath ID 3153
% Compile twice: lualatex 3153.tex
\documentclass[11pt]{article}
\usepackage{fontspec}
\setmainfont{Latin Modern Roman}
\usepackage{amsmath,amssymb,mathtools,unicode-math}
\setmathfont{Latin Modern Math}
\usepackage{xurl,hyperref}
\hypersetup{hidelinks,pdftitle={UnsolvedMath Problem 3153}}
\newcommand{\sref}[2]{\hyperlink{source:#1:#2}{[S#2]}}
\newcommand{\cataloguescope}[1]{\paragraph{Mathematical question.}#1}
\begin{document}
% UnsolvedMath ID 3153; source catalogue code OPG-826
\setcounter{section}{3152}
\section{Complete bipartite subgraphs of perfect graphs}\label{3153}
{\small\sffamily UnsolvedMath ID 3153 · Graph Theory}
\subsection{Definitions and mathematical statement}

\paragraph{Shared notation.}$\mathbb N=\{1,2,\ldots\}$, and $\mathbb Z,\mathbb Q,\mathbb R,\mathbb C$ denote the integers, rationals, reals and complex numbers. Any other conventions are specified locally.


A finite simple graph $G$ is perfect if every induced subgraph $H$ satisfies $\chi(H)=\omega(H)$, where these are chromatic and clique numbers. A complete bipartite subgraph with parts $A,B$ consists of disjoint vertex sets for which every vertex of $A$ is adjacent to every vertex of $B$; edges inside the parts are irrelevant. Let $\overline G$ denote the complement.
\paragraph{Statement recorded in the supplied source.}
Is it true that for every $\varepsilon>0$ there is $n_0(\varepsilon)$ such that, for every perfect graph $G$ on $n\ge n_0(\varepsilon)$ vertices, one of $G,\overline G$ has a complete bipartite subgraph with $|A|,|B|\ge n^{1-\varepsilon}$? This is the uniform asymptotic meaning of the source's $n^{1-o(1)}$ bound. \sref{3153}{2}
\subsection{Short English statement}
Does every large perfect graph or its complement contain a balanced complete bipartite subgraph of nearly linear size?
\subsection{Sources}
\begin{description}
\item[\hypertarget{source:3153:1}{S1}] ULAM.AI and the UnsolvedMath Contributors, UnsolvedMath: A Curated Collection of Open Mathematics Problems, record 3153 (OPG-826). CC BY 4.0. \url{https://huggingface.co/datasets/ulamai/UnsolvedMath}.
\item[\hypertarget{source:3153:2}{S2}] Open Problem Garden, Complete bipartite subgraphs of perfect graphs, original node 826, statement and mathematical discussion. \url{https://www.openproblemgarden.org/op/complete_bipartite_subgraphs_of_perfect_graphs}.
\end{description}
\label{end:3153}
\end{document}
